\documentclass[a4paper,10pt]{article}

\usepackage[utf8]{inputenc}
\usepackage[T1]{fontenc}
\usepackage[turkish]{babel}
\usepackage{geometry}
\usepackage{booktabs}
\usepackage{array}
\usepackage{tabularx}
\usepackage{longtable}
\usepackage{enumitem}
\usepackage{listings}
\usepackage{xcolor}
\usepackage[hidelinks]{hyperref}
\usepackage{microtype}
\usepackage{amsmath}
\usepackage{titlesec}

\AtBeginDocument{\shorthandoff{=}}

\geometry{
    a4paper,
    top=1.8cm,
    bottom=1.8cm,
    left=1.8cm,
    right=1.8cm
}

\setlength{\parindent}{0pt}
\setlength{\parskip}{5pt}
\renewcommand{\arraystretch}{1.18}

\titleformat{\section}{\large\bfseries}{\thesection.}{0.5em}{}
\titleformat{\subsection}{\normalsize\bfseries}{\thesubsection.}{0.5em}{}

\lstdefinestyle{cstyle}{
    language={[ANSI]C},
    basicstyle=\ttfamily\small,
    keywordstyle=\bfseries,
    commentstyle=\itshape,
    numbers=left,
    numberstyle=\tiny,
    numbersep=7pt,
    frame=single,
    breaklines=true,
    showstringspaces=false,
    tabsize=4
}

\newcommand{\hex}[1]{\texttt{0x#1}}
\newcommand{\code}[1]{\texttt{#1}}

\title{\textbf{DEOS Traction Node ICD}\\
\large v1.0 (PDF Revizyonu)}
\author{DEOS -- Dokuz Eylül Otonom Sistemler}
\date{28 Eylül 2026}

\begin{document}
\maketitle

\section{Traction Node Commands}

Tahrik (Traction) sistemi \code{DEOS\_NODE\_TRACTION} (\hex{10}) adresindedir. 

\subsection{Gerçek Zamanlı Kontrol Komutları (Real-Time RX)}

Aşağıdaki komutlar Traction ECU'nun ağdan anlık olarak alması gereken komutları tanımlar.

\begin{center}
\begin{tabularx}{\textwidth}{>{\raggedright\arraybackslash}p{3.8cm} >{\raggedright\arraybackslash}p{3cm} >{\raggedright\arraybackslash}p{2.2cm} >{\raggedright\arraybackslash}p{2.2cm} X}
\toprule
\textbf{Değişken Adı} & \textbf{Sınıf \& Servis} & \textbf{Komut ID} & \textbf{Ölçek \& Birim} & \textbf{Fonksiyonel Görev} \\
\midrule
\code{target\_speed} & COMMAND (0x1) TRACTION (0x01) & 0x01 (SET\_SPEED) & 0.01 m/s & micro-ROS veya Textual referans hız komutu. \\
\code{target\_accel} & COMMAND (0x1) TRACTION (0x01) & 0x01 (Genişletilmiş) & 0.01 m/s² & İleri besleme girdisi; dişli boşluğu (backlash) darbe önleyici. \\
\code{target\_gear} & COMMAND (0x1) TRACTION (0x01) & 0x04 (SET\_GEAR) & Enum (0:N, 1:D, 2:R, 3:P) & Vites seçimi. Hız sıfırlanmadan ters yöne geçiş kilitlenir. \\
\code{dynamic\_torque\_limit} & COMMAND (0x1) TRACTION (0x01) & 0x05 (SET\_TORQUE) & 0.01 \% & Sürüş moduna göre anlık motor tork tavanı sınırlandırması. \\
\code{emergency\_stop\_cmd} & SAFETY (0x0) SYSTEM (0x00) & 0x05 (E-STOP) & 0:Normal 1:E-STOP & Donanımsal PWM'i \%0'a çeker, motor akımını keser. \\
\code{brake\_system\_active} & STATUS (0x2) BRAKE (0x03) & 0x02 (GET\_STATUS) & 0:Boşta 1:Aktif & Fren basılıyken gaz referansını derhal sıfırlar. (Fren-Gaz Kilidi) \\
\bottomrule
\end{tabularx}
\end{center}

\subsection{BMS Dinamik Güç Koruması (Çapraz Telemetri)}
Traction node'u, BMS'den gelen şu yayınları (0x20) dinleyerek kendi sınırlarını dinamik olarak ayarlar:
\begin{itemize}
    \item \code{bms\_battery\_soc}: SoC > \%95 olduğunda rejeneratif freni derhal kapatır.
    \item \code{bms\_max\_discharge\_current}: Çekiş (tahrik) akım tavanı.
    \item \code{bms\_max\_charge\_current}: Rejeneratif şarj akımı tavanı.
\end{itemize}

\section{Traction Konfigürasyon Parametreleri}

Aşağıdaki parametreler \code{CLASS\_CONFIG (0x6)}, \code{SERVICE\_TRACTION (0x01)} üzerinden konfigüre edilir.

\begin{center}
\begin{tabularx}{\textwidth}{>{\raggedright\arraybackslash}p{3.8cm} >{\raggedright\arraybackslash}p{3cm} >{\raggedright\arraybackslash}p{2.2cm} >{\raggedright\arraybackslash}p{2.2cm} X}
\toprule
\textbf{Parametre} & \textbf{ID} & \textbf{Tip} & \textbf{Ölçek} & \textbf{Açıklama} \\
\midrule
PID\_KP & \hex{0001} & uint16\_t & 0.01 & Hız takip döngüsü oransal kazancı. \\
PID\_KI & \hex{0002} & uint16\_t & 0.01 & Hız takip döngüsü integral kazancı. \\
PID\_KD & \hex{0003} & uint16\_t & 0.01 & Hız takip döngüsü türev kazancı. \\
MAX\_CURRENT & \hex{0010} & uint16\_t & 0.1 A & Maksimum motor faz akım tavanı. \\
MAX\_SPEED & \hex{0011} & uint16\_t & 0.01 m/s & İzin verilen azami araç sürati. \\
ACCEL\_LIMIT & \hex{0012} & uint16\_t & 0.01 m/s² & Azami ivmelenme rampa sınırı (Slew-rate). \\
CMD\_TIMEOUT\_MS & \hex{0013} & uint16\_t & 1 ms & Otonom komut kesilme emniyet süresi. \\
POLE\_PAIRS & \hex{0014} & uint8\_t & 1 & BLDC motor kutup çifti sayısı. \\
GEAR\_RATIO & \hex{0015} & uint16\_t & 0.01 & Aktarma organı redüktör tahvil oranı. \\
WHEEL\_CIRCUM\_MM & \hex{0020} & uint16\_t & 1 mm & Tekerlek dinamik yuvarlanma çevresi. \\
REGEN\_CUTOFF\_SPD & \hex{0022} & uint16\_t & 0.01 m/s & Rejen frenin kapatılacağı alt hız eşiği. \\
\bottomrule
\end{tabularx}
\end{center}

\section{Traction Kalibrasyon ve Donanım İşlemleri}
Traction kalibrasyonu \code{DEOS\_SERVICE\_CALIBRATION (0x05)} servisi üzerinden yürütülür:
\begin{itemize}
    \item \textbf{Gaz Sıfır Gerilim Ofseti:} STM32 gaz PWM'ini \%0 yapar; RC filtre ve op-amp tampon çıkışındaki analog gerilim okunarak sıfır ofseti EEPROM'a kaydedilir (High-Pedal kilidini önler).
    \item \textbf{Dinamik Yuvarlanma Çevresi:} 10 metrelik test odometrisi.
    \item \textbf{Donanım:} Optokuplörler (6N137, PC817), 159.15 Hz RC filtre ve MCP6002 tampon ile 0–5V DC dönüştürücü kullanılır.
\end{itemize}

\end{document}
